\begin{textsample}{POS Dim 1 – df – Score 33.00 – df/df\_em\_13.txt}  \label{ex:f1_pos_096}
4.
AVALIAÇÃO A avaliação \textbf{configura} -se como \textbf{um} dos pilares fundamentais do trabalho pedagógico da escola e da sala de aula e perpassa por todo o processo, estando presente desde o início até o final.
Avaliação e aprendizagem caminham juntas no processo educativo.
Dessa forma, não há avaliação sem aprendizagem, assim como não há aprendizagem sem avaliação.
Faz -se necessário compreender \textbf{que} > [... ] avaliar é \textbf{um} processo natural \textbf{que} faz parte do repertório das ações dos seres humanos e está indissociavelmente ligado à busca de todo tipo de ação \textbf{que} provoque mudanças, mudanças essas \textbf{que} indicam respostas concretas ao desejo por “ algo melhor ”.
O ato de avaliar tem como foco a construção dos melhores resultados possíveis.
O ser humano permanentemente avalia em busca de \textbf{uma} melhor qualidade de vida, da construção de si mesmo e do melhor modo de ser e \textbf{viver} ( \textbf{SANTOS} ; GONTIJO, 2018, p. 31 ).
Villas Boas ( 2014 ) utiliza a expressão “ avaliação para aprendizagem ”, conotação semântica \textbf{que} condiz com a emergente constituição de \textbf{uma} nova concepção de avaliação, \textbf{baseada} em processos permanentes de reflexão e comprometida com a aprendizagem e o desenvolvimento humano dos estudantes. \textbf{Ademais}, “ a avaliação para aprendizagem tem a conotação de movimento, de busca pela aprendizagem, pelo professor e pelos alunos, \textbf{enquanto} avaliação das aprendizagens se ocupa do processo já ocorrido ” ( VILLAS BOAS, 2014, p. 68 ).
Em função da proximidade dos conceitos, a avaliação para aprendizagem também pode ser tratada sob a denominação de avaliação formativa.
Esta corrobora com o princípio de \textbf{que} a prática avaliativa deve tornar -se auxiliar à aprendizagem.
A esperança de dispor da avaliação a serviço da aprendizagem e a certeza de \textbf{que} isso possui teor \textbf{legítimo} nas situações pedagógicas indicam \textbf{que} a avaliação formativa \textbf{configura} -se como \textbf{um} \textbf{horizonte} lógico das práticas avaliativas no contexto escolar ( HADJI, 2001 ).
A avaliação formativa constitui -se, portanto, em \textbf{um} processo complexo e cujo detalhamento dar -se-á na relação direta entre professores e estudantes, ou seja, deve ser dimensionada e modulada para cada realidade escolar, não havendo \textbf{uma} fórmula a ser aplicada, mas, sim, o desenvolvimento de \textbf{um} processo \textbf{que} é parte da própria aprendizagem do estudante do fazer pedagógico do professor.
Para Villas Boas : > [... ] a avaliação cumpre, também, função formativa pela qual os professores analisam, de maneira frequente e interativa, o progresso dos alunos, para identificar o \textbf{que} eles aprenderam e o \textbf{que} ainda não aprenderam, para \textbf{que} venham a aprender, e para \textbf{que} reorganizem o trabalho pedagógico.
Essa avaliação requer \textbf{que} se considerem as diferenças dos alunos, se adapte o trabalho às necessidades de cada \textbf{um} e se dê tratamento adequado de seus resultados. ( VILLAS BOAS, 2006 ).
A avaliação formativa possui \textbf{um} papel \textbf{central} no \textbf{que} se refere à organização do Currículo em Movimento para o Ensino Médio, \textbf{embasado} na repartição do tempo escolar em duas partes indissociáveis — Formação Geral Básica e Itinerários Formativos, ambas articuladas pelo Projeto de Vida.
Nesse sentido, a avaliação formativa atua como processo \textbf{metodológico} essencial para \textbf{que} sejam atingidos os objetivos de aprendizagem propostos no currículo.
Isso visa manter os pressupostos \textbf{que} levaram à elaboração deste documento e adequá -los, naquilo \textbf{que} for necessário, à nova legislação vigente.
Nos últimos anos, o Ensino Médio tem passado por \textbf{inúmeras} mudanças, tendo como marcos legais \textbf{norteadores} a Lei nº 13.415/2017, as DCNEM/2018, a BNCC/2018, os Referenciais Curriculares para a elaboração dos Itinerários Formativos, de 2018, e a Nota Técnica nº 2/2019 do Conselho de Educação do Distrito Federal ( CEDF ).
Dentre os documentos \textbf{norteadores} já mencionados, cabe destacar os parágrafos acrescidos ao art. 35 -A da LDB, pela Lei nº 13.415/2017 : > § 7º Os currículos do ensino médio deverão considerar a formação integral do aluno, de maneira a adotar \textbf{um} trabalho voltado para a construção de seu projeto de vida e para sua formação nos aspectos físicos, cognitivos e socioemocionais. > > § 8º Os conteúdos, as metodologias e as formas de avaliação processual e formativa serão organizados nas redes de ensino por meio de atividades \textbf{teóricas} e práticas, \textbf{provas} orais e escritas, seminários, projetos e atividades on-line, de tal forma \textbf{que} ao final do ensino médio o \textbf{educando} demonstre : > > I - domínio dos princípios científicos e tecnológicos \textbf{que} presidem a produção \textbf{moderna} ; > > II - conhecimento das formas \textbf{contemporâneas} de linguagem ( BRASIL, 2017 ).
A educação escolar orienta -se intencionalmente por metas \textbf{que} objetivam acompanhar todo o processo de ensino e aprendizagem.
Tais intenções da ação educativa só adquirem sentido quando levadas em consideração a natureza social e a função socializadora da educação escolar, as quais têm como razão primordial a promoção do desenvolvimento humano.
Incluem -se nesse processo os procedimentos didáticos assumidos pelo professor e a avaliação como ferramenta fundamental na aprendizagem ( DARSIE, 1996 ).
Na perspectiva de \textbf{uma} concepção de currículo como construto social, fica evidente \textbf{que} o papel da avaliação é \textbf{uma} ação \textbf{que} se projeta \textbf{focada} sempre no futuro, comprometida com o sucesso de todos, superando o modelo de avaliação classificatório, autoritário e excludente por \textbf{um} modelo favorecedor da aprendizagem.
Hoffmann ( 2008 ) clarifica tal ação ao realçar \textbf{que} as práticas avaliativas tradicionais detêm -se nos aspectos comprobatórios de \textbf{uma} etapa vencida pelo estudante, por meio de resultados obtidos no processo educativo, tecidos por considerações atitudinais \textbf{que} justificam ou explicam o alcance de tais resultados em determinado espaço de tempo.
Assim, a avaliação estrutura -se no passado, relatando e explicando o presente.
O modelo avaliativo tradicional tem como alvo apenas o estudante.
A implementação do Novo Ensino Médio no Distrito Federal sinaliza \textbf{uma} necessária ressignificação da escola como espaço de convivência e de relações sociais, de aprendizagens essenciais e de oferta de trajetórias diversificadas.
Nesse sentido, os novos tempos e contornos legais apontam para \textbf{uma} proposta avaliativa transformadora e possível, \textbf{demandando} \textbf{que} o estudante tome consciência de seu processo de aprendizagem, sendo protagonista de seu próprio desenvolvimento pessoal e educativo.
Essa prática favorece substancialmente o exercício da metacognição[^6 ] ou da meta-aprendizagem, transformando -a em instrumento ativo e real de aprendizagem ( DARSIE, 1996 ).
Quanto ao envolvimento dos estudantes no processo avaliativo, consolida -se a possibilidade de eles se tornarem parceiros dessa importante \textbf{tarefa}.
Assim, são convidados a participarem da definição dos critérios da avaliação e a aplicá -los em seus trabalhos.
Deve -se criar a cultura da avaliação desvinculada da nota e da promoção ou reprovação, reforçando a ideia de \textbf{que} todos são capazes de aprender.
Naturalmente, nesse sentido, constrói -se a confiança e a segurança.
O envolvimento dos estudantes na \textbf{tarefa} de registrar os resultados lhes possibilita \textbf{um} acompanhamento de seu desempenho por meio de \textbf{uma} autoavaliação contínua.
Essa concepção requer a avaliação dos estudantes em sua plenitude, incluindo a formação cidadã, cujo objetivo é a inserção social crítica.
O envolvimento dos estudantes no processo avaliativo lhes permite a \textbf{partilha} de informações com outros sobre seu progresso ( VILLAS BOAS, 2007 ).
Paralelamente e considerando as alterações ocorridas na LDB pela Lei nº 13.415/2017, evidencia -se \textbf{que} houve atualização das DCNEM, apontando princípios e \textbf{fundamentos} para a organização das políticas públicas educacionais.
Destaca -se, sobremaneira, como \textbf{um} dos princípios apontados, o Projeto de Vida, \textbf{que} é considerado “ como estratégia de reflexão sobre a trajetória escolar na construção das dimensões pessoal, cidadã e profissional do estudante ” ( BRASIL, 2018b ).
Na atualização das DCNEM, o Projeto de Vida apresenta -se como \textbf{uma} real estratégia pedagógica, com o objetivo de promoção do autoconhecimento do estudante e de sua dimensão cidadã, orientando o planejamento da carreira profissional almejada, ao considerar seus interesses, talentos, desejos e potencialidades ( SILVA, CARVALHO, 2019 ).
Na concepção de avaliação formativa, todos avaliam todos.
Nessa perspectiva, acredita -se \textbf{que} a função formativa oportuniza a promoção das aprendizagens de todos por meio da autoavaliação e do feedback ( devolutiva/retorno ).
Professor e estudante compõem \textbf{dialeticamente} \textbf{um} movimento no qual todos os \textbf{atores} do processo educativo tomam como elemento valioso o diálogo \textbf{que} \textbf{ora} se estabelece.
Tal concepção \textbf{exige} \textbf{um} processo formativo \textbf{que} favoreça, permanentemente, o desenvolvimento profissional docente em relação aos aspectos \textbf{teóricos} e críticos, bem como \textbf{exige} a elaboração coletiva dos objetivos do trabalho pedagógico.
A avaliação formativa é \textbf{um} processo permanente em construção, tem \textbf{um} caráter processual e contínuo e adquire o caráter concomitante de orientação e reorientação da aprendizagem.
Por fim, os estudantes de Ensino Médio, para alcançarem determinado aprendizado, deparam -se com \textbf{um} processo complexo, de cunho pessoal, para o qual os professores e a escola podem contribuir ao permitir a comunicabilidade, ao situá -los em seu grupo e ao oportunizar \textbf{que} aprendam a respeitar e a fazer -se respeitar.
Para tal, é salutar \textbf{que} professores e escola criem situações em \textbf{que} o estudante deva sentir -se desafiado ou \textbf{instigado} a participar e questionado pelo jogo do conhecimento, adquirindo o \textbf{espírito} de pesquisa e o desenvolvimento de capacidades de \textbf{raciocínio} e de autonomia ( BRASIL, 1999 ). [ ^6 ] : Entendemos a metacognição como sendo a capacidade de avaliar, regular e organizar os processos mentais sobre os objetos cognitivos estudados para \textbf{que} se possa atingir \textbf{uma} meta ou \textbf{um} objetivo.
Podemos, ainda, traduzir a metacognição como sendo o pensar sobre o pensar.

% matched lemmas: ademais, ator, basear, central, configurar, contemporâneo, demandar, dialeticamente, educar, embasar, enquanto, espírito, exigir, focar, fundamento, horizonte, instigar, inúmero, legítimo, metodológico, moderno, norteador, ora, partilha, prova, que, raciocínio, santo, tarefa, teórico, um, viver
\end{textsample}
